\documentclass{article}
\usepackage{algorithm, algorithmic, amsfonts, amsmath}

\begin{document}
\begin{algorithm}[!h]
\renewcommand{\algorithmicrequire}{\textbf{Input:}} % show \REQUIRE as Input
\renewcommand{\algorithmicensure}{\textbf{Output:}} % show \ENSURE as Output
	\caption{PULSNAR (class-proportion estimation under SNAR)} % algorithm title
	\begin{algorithmic}[1] % [1] numbers every line
	    \REQUIRE
	    Data $X$; label indicator $Y\in\{0,1\}$; classifier (e.g. XGBoost);
	    clustering method (e.g. GMM); maximum clusters $K_{\max}$;
	    number of iterations $n_{it}$; calibration method
	    \ENSURE
	    Estimated positive fraction $\hat{\alpha}$; calibrated probabilities
	    \STATE \textbf{Setting:} SNAR --- positives are \emph{selected non-at-random}, so the
	    labeled positives are not representative of the positive class % a weaker assumption than SAR: different positive subgroups are labeled at different rates
	    \IF{SCAR is assumed}
	    \STATE \textbf{PULSCAR branch}
	    \FOR{$itr=1$ \TO $n_{it}$}
	    \STATE Shuffle $(X,Y)$; train the classifier and estimate $\hat{\alpha}$ from the
	    predicted-probability histogram % a single estimate has high variance, so several rounds are pooled
	    \ENDFOR
	    \ELSE
	    \STATE \textbf{PULSNAR branch}
	    \STATE Find important features using all data % clustering on important features rather than all of them reduces dimensional noise
	    \STATE Divide the data into positive and unlabeled sets
	    \STATE Cluster the \emph{labeled positives} into $k$ clusters;
	    if $k$ is not given, select it by BIC over $2..K_{\max}$ % the central assumption: positives form subgroups, and SNAR shows up as differing labeling rates across them
	    \FOR{each positive cluster $j=1,\dots,k$}
	    \STATE Treat cluster $j$ against the unlabeled set as a \emph{sub-problem} satisfying SCAR
	    \STATE Estimate the sub-proportion $\hat{\alpha}_{j}$ by the PULSCAR procedure % SCAR holds approximately inside a subgroup, so a SCAR method applies there
	    \ENDFOR
	    \STATE $\hat{\alpha}\leftarrow \sum_{j=1}^{k}\hat{\alpha}_{j}$ % summing the subgroup proportions recovers the overall positive proportion
	    \ENDIF
	    \STATE \textbf{Probability calibration} (optional): map the classifier's scores using
	    $\hat{\alpha}$, e.g. by isotonic regression % the output probabilities of a PU classifier are systematically low and need calibrating against the estimated alpha
		\RETURN $\hat{\alpha}$, calibrated probabilities, performance metrics
	\end{algorithmic}
\end{algorithm}
% Kumar & Lambert, PeerJ Computer Science 2024, "PULSNAR --- Positive unlabeled learning
% selected not at random". Compiled from the main flow of the authors' released implementation
% (scar_data_processing / snar_data_processing in PULSNAR/PULSNAR.py,
% AlphaEstimate.py, ProbabilityCalibration.py).
% Difference between SNAR and SAR: SAR assumes the labeling probability is a function e(x) of the instance, whereas
% SNAR goes further and holds that the positive class is composed of heterogeneous subgroups labeled at different
% rates. Medicine is the typical case, where subtypes of one disease differ widely in how often they are recorded.
% The key step is divide and conquer: cluster the positives into k subgroups, within each of which SCAR holds
% approximately, apply the SCAR proportion estimator (PULSCAR) there, and sum the k sub-proportions.
% In this project PULSNAR is recorded only as provenance for the SNAR direction and enters no numerical result.
\end{document}
